\begingroup
\renewcommand{\arraystretch}{1.45}
\SourceLeaf{292}
\[
\begin{aligned}
W={}&\frac{A\cos(\tfrac12a+k)}{\sin\tfrac12(a-b)\sin\tfrac12(a-c)\sin\tfrac12(a-d)\mathinner{.\,.\,.\,.}\sin\tfrac12(a-t)}\\[5pt]
&+\frac{B\cos(\tfrac12b+k)}{\sin\tfrac12(b-a)\sin\tfrac12(b-c)\sin\tfrac12(b-d)\mathinner{.\,.\,.\,.}\sin\tfrac12(b-t)}\\[5pt]
&+\frac{C\cos(\tfrac12c+k)}{\sin\tfrac12(c-a)\sin\tfrac12(c-b)\sin\tfrac12(c-d)\mathinner{.\,.\,.\,.}\sin\tfrac12(c-t)}\\[5pt]
&+\frac{D\cos(\tfrac12d+k)}{\sin\tfrac12(d-a)\sin\tfrac12(d-b)\sin\tfrac12(d-c)\mathinner{.\,.\,.\,.}\sin\tfrac12(d-t)}\\[5pt]
&+\text{etc.}\\[5pt]
&+\frac{T\cos(\tfrac12t+k)}{\sin\tfrac12(t-a)\sin\tfrac12(t-b)\sin\tfrac12(t-c)\sin\tfrac12(t-d)\mathinner{.\,.\,.\,.}}
\end{aligned}
\]
Quodsi iam quantitatem $k$, quae prorsus arbitraria est, $=-\tfrac12t$ ponimus,
summam autem $a+b+c+d+\mathinner{.\,.\,.}=s$ (exclusa $t$), prodit, multiplicando per
$\sin\tfrac12(t-a)\sin\tfrac12(t-b)\sin\tfrac12(t-c)\sin\tfrac12(t-d)\mathinner{.\,.\,.\,.}$
\[
\begin{aligned}
T={}&\frac{\sin\tfrac12(t-b)\sin\tfrac12(t-c)\sin\tfrac12(t-d)\mathinner{.\,.\,.\,.}}{\sin\tfrac12(a-b)\sin\tfrac12(a-c)\sin\tfrac12(a-d)\mathinner{.\,.\,.\,.}}A\cos\tfrac12(t-a)\\[5pt]
&+\frac{\sin\tfrac12(t-a)\sin\tfrac12(t-c)\sin\tfrac12(t-d)\mathinner{.\,.\,.\,.}}{\sin\tfrac12(b-a)\sin\tfrac12(b-c)\sin\tfrac12(b-d)\mathinner{.\,.\,.\,.}}B\cos\tfrac12(t-b)\\[5pt]
&+\frac{\sin\tfrac12(t-a)\sin\tfrac12(t-b)\sin\tfrac12(t-d)\mathinner{.\,.\,.\,.}}{\sin\tfrac12(c-a)\sin\tfrac12(c-b)\sin\tfrac12(c-d)\mathinner{.\,.\,.\,.}}C\cos\tfrac12(t-c)\\[5pt]
&+\frac{\sin\tfrac12(t-a)\sin\tfrac12(t-b)\sin\tfrac12(t-c)\mathinner{.\,.\,.\,.}}{\sin\tfrac12(d-a)\sin\tfrac12(d-b)\sin\tfrac12(d-c)\mathinner{.\,.\,.\,.}}D\cos\tfrac12(t-d)\\[5pt]
&+\text{etc.}\\[5pt]
&+2^{2m-1}i^{2m}\sin\tfrac12(t-a)\sin\tfrac12(t-b)\sin\tfrac12(t-c)\sin\tfrac12(t-d)\mathinner{.\,.\,.}\\[5pt]
&\qquad\times\{\alpha^m\cos\tfrac12s+\beta^m\sin\tfrac12s\}
\end{aligned}
\]
Circa hanc formulam quasdam adhuc annotationes adiicimus

I. Est $i^{2m}=+1$, vel $-1$, prout $m$ par est vel impar.

II. Quum formula indefinite pro quovis valore ipsius $t$ valeat, necessario,
evolutis sinuum productis in sinus et cosinus arcuum multiplicium, cum formula
pro $X$, si pro $x$ scribitur $t$, identica erit, i. e. cum
\[
\begin{aligned}
&\alpha+\alpha'\cos t+\alpha''\cos2t+\text{etc.}\\[5pt]
&\quad+\beta'\sin t+\beta''\sin2t+\text{etc.}
\end{aligned}
\]
III. Determinantur itaque omnes coëfficientes $\alpha,\alpha',\beta',\alpha'',\beta''$ etc. per quantitates cognitas, atque unicam incognitam
$\alpha^m\cos\tfrac12s+\beta^m\sin\tfrac12s$, quamobrem, si
qua inter illos relatio insuper datur, omnes ex asse facile poterunt assignari.
\par\SourceLeaf{293}
IV. Quodsi autem tantummodo expressio eiusdem formae ut
\[
\begin{aligned}
&\alpha+\alpha'\cos t+\alpha''\cos2t+\text{etc.}\\[5pt]
&\quad+\beta'\sin t+\beta''\sin2t+\text{etc.}
\end{aligned}
\]
quaeritur (sive cum hac identica sit, sive non), per quam valores propositi
$A,B,C,D$ etc. repraesententur, si pro $t$ resp. $a,b,c,d$ etc. substituitur: quantitas
$\alpha^m\cos\tfrac12s+\beta^m\sin\tfrac12s$ ex incognita \emph{arbitraria} evadit, quam si placet etiam $=0$
statuere licebit; reveraque formula nostra pro $T$ ita comparata est, ut absque
calculo sponte pateat, hacce eam proprietate praeditam esse.

\InterpolationArticle{12}
Operae pretium erit, transformationem ei analogam, quam in art. 4 explicavimus,
etiam ad formulas art. 10 et art. praec., si ibi $\alpha^m\cos\tfrac12s+\beta^m\sin\tfrac12s=0$
ponitur, applicare, quo pacto hos duos casus diversos sub formula unica comprehendere licebit.
Crescat itaque multitudo quantitatum $a,b,c,d\mathinner{.\,.\,.}$ successive
ab una $a$, ad duas $a,b$, deinde ad tres $a,b,c$ etc.; evolvatur alternis vicibus
per artt. 10 et 11 valor, quem $t$ pro singulis his suppositionibus nanciscitur, tandem
quisque valor a proxime sequente subtrahatur, sequenti modo:
\[
\begin{aligned}
&A\\[5pt]
&\frac{\sin\tfrac12(t-b)}{\sin\tfrac12(a-b)}A\cos\tfrac12(t-a)\\[5pt]
&+\frac{\sin\tfrac12(t-a)}{\sin\tfrac12(b-a)}B\cos\tfrac12(t-b)\\[5pt]
&\frac{\sin\tfrac12(t-b)\sin\tfrac12(t-c)}{\sin\tfrac12(a-b)\sin\tfrac12(a-c)}A\\[5pt]
&+\frac{\sin\tfrac12(t-a)\sin\tfrac12(t-c)}{\sin\tfrac12(b-a)\sin\tfrac12(b-c)}B\\[5pt]
&+\frac{\sin\tfrac12(t-a)\sin\tfrac12(t-b)}{\sin\tfrac12(c-a)\sin\tfrac12(c-b)}C\\[5pt]
&\text{etc.}
\end{aligned}
\]
Hic differentia inter valorem primum et secundum fit
\[
=\sin\tfrac12(t-a)\cos\tfrac12(t-b)
\left\{\frac{A}{\sin\tfrac12(a-b)}+\frac{B}{\sin\tfrac12(b-a)}\right\}
\]
Differentia inter valorem secundum et tertium
\[
\begin{aligned}
={}&\sin\tfrac12(t-a)\sin\tfrac12(t-b)\left\{
\frac{A\cos\tfrac12(a-c)}{\sin\tfrac12(a-b)\sin\tfrac12(a-c)}\right.\\[5pt]&\qquad\left.
+\frac{B\cos\tfrac12(b-c)}{\sin\tfrac12(b-a)\sin\tfrac12(b-c)}\right.\\[5pt]&\qquad\left.
+\frac{C}{\sin\tfrac12(c-a)\sin\tfrac12(c-b)}\right\}
\end{aligned}
\]
Quae operatio si ultra continuatur, facile emergit lex generalis. Scilicet statuendo
\SourceLeaf{294}
\[
\begin{aligned}
A'={}&\frac{A}{\sin\tfrac12(a-b)}+\frac{B}{\sin\tfrac12(b-a)}\\[5pt]
A''={}&\frac{A\cos\tfrac12(a-c)}{\sin\tfrac12(a-b)\sin\tfrac12(a-c)}\\[5pt]
&+\frac{B\cos\tfrac12(b-c)}{\sin\tfrac12(b-a)\sin\tfrac12(b-c)}\\[5pt]
&+\frac{C}{\sin\tfrac12(c-a)\sin\tfrac12(c-b)}\\[5pt]
A'''={}&\frac{A}{\sin\tfrac12(a-b)\sin\tfrac12(a-c)\sin\tfrac12(a-d)}\\[5pt]
&+\frac{B}{\sin\tfrac12(b-a)\sin\tfrac12(b-c)\sin\tfrac12(b-d)}\\[5pt]
&+\frac{C}{\sin\tfrac12(c-a)\sin\tfrac12(c-b)\sin\tfrac12(c-d)}\\[5pt]
&+\frac{D}{\sin\tfrac12(d-a)\sin\tfrac12(d-b)\sin\tfrac12(d-c)}\\[5pt]
\end{aligned}\]
\[\begin{aligned}
A''''={}&\frac{A\cos\tfrac12(a-e)}{\sin\tfrac12(a-b)\sin\tfrac12(a-c)\sin\tfrac12(a-d)\sin\tfrac12(a-e)}\\[5pt]
&+\frac{B\cos\tfrac12(b-e)}{\sin\tfrac12(b-a)\sin\tfrac12(b-c)\sin\tfrac12(b-d)\sin\tfrac12(b-e)}\\[5pt]
&+\frac{C\cos\tfrac12(c-e)}{\sin\tfrac12(c-a)\sin\tfrac12(c-b)\sin\tfrac12(c-d)\sin\tfrac12(c-e)}\\[5pt]
&+\frac{D\cos\tfrac12(d-e)}{\sin\tfrac12(d-a)\sin\tfrac12(d-b)\sin\tfrac12(d-c)\sin\tfrac12(d-e)}\\[5pt]
&+\frac{E}{\sin\tfrac12(e-a)\sin\tfrac12(e-b)\sin\tfrac12(e-c)\sin\tfrac12(e-d)}
\end{aligned}
\]
etc. fiet
\[
\begin{aligned}
T={}&A+A'\sin\tfrac12(t-a)\cos\tfrac12(t-b)\\[5pt]
&+A''\sin\tfrac12(t-a)\sin\tfrac12(t-b)\\[5pt]
&+A'''\sin\tfrac12(t-a)\sin\tfrac12(t-b)\sin\tfrac12(t-c)\cos\tfrac12(t-d)\\[5pt]
&+A''''\sin\tfrac12(t-a)\sin\tfrac12(t-b)\sin\tfrac12(t-c)\sin\tfrac12(t-d)\\[5pt]
&+A^{\mathrm{V}}\sin\tfrac12(t-a)\sin\tfrac12(t-b)\sin\tfrac12(t-c)\\[5pt]&\qquad \sin\tfrac12(t-d)\sin\tfrac12(t-e)\cos\tfrac12(t-f)\\[5pt]
&+A^{\mathrm{VI}}\sin\tfrac12(t-a)\sin\tfrac12(t-b)\sin\tfrac12(t-c)\\[5pt]&\qquad \sin\tfrac12(t-d)\sin\tfrac12(t-e)\sin\tfrac12(t-f)\\[5pt]
&+\text{etc.}
\end{aligned}
\]
quae progressio ad totidem terminos continuanda est, quot valores functionis $X$
dati sunt.

Coëfficientes $A',A'',A''',A''''$ etc. etiam per algorithmum sequentem computari possunt.
Designetur per $B',B'',B''',B''''$ etc. id quod illi resp. fiunt, si
\[
a,\quad A,\quad b,\quad B,\quad c,\quad C,\quad d,\quad D\quad\text{etc.}
\]
resp. mutantur in
\[
b,\quad B,\quad c,\quad C,\quad d,\quad D,\quad e,\quad E\quad\text{etc.}
\]
Porro transeant
\[\begin{gathered}
A',A'',A''',A''''\text{ etc.}\quad\text{in}\quad C',C'',C''',C''''\text{ etc.}\quad
\\[5pt]\text{vel in}\quad D',D'',D''',D''''\text{ etc. etc.}
\end{gathered}\]
si
\SourceLeaf{295}
\[
a,\ A,\ b,\ B,\ c,\ C,\ d,\ D\quad\text{etc.}
\]
resp. mutantur in
\[
c,\ C,\ d,\ D,\ e,\ E,\ f,\ F\quad\text{etc.}
\]
vel in
\[
d,\ D,\ e,\ E,\ f,\ F,\ g,\ G\quad\text{etc.}
\]
\hspace*{2em}etc.

Tunc erit
\[\begin{gathered}
A'=\frac{A-B}{\sin\tfrac12(a-b)},\\[5pt]
B'=\frac{B-C}{\sin\tfrac12(b-c)},\\[5pt]
C'=\frac{C-D}{\sin\tfrac12(c-d)}\quad\text{etc.}
\end{gathered}\]
\[\begin{gathered}
A''=\frac{A'\cos\tfrac12(a-c)-B'}{\sin\tfrac12(a-c)},\\[5pt]
B''=\frac{B'\cos\tfrac12(b-d)-C'}{\sin\tfrac12(b-d)}\quad\text{etc.}
\end{gathered}\]
\[\begin{gathered}
A'''=\frac{A''\cos\tfrac12(a-c)+A'\sin\tfrac12(a-c)-B''}{\sin\tfrac12(a-d)},\\[5pt]
B'''=\frac{B''\cos\tfrac12(b-d)+B'\sin\tfrac12(b-d)-C''}{\sin\tfrac12(b-e)}\quad\text{etc.}
\end{gathered}\]
\[\begin{gathered}
A''''=\frac{A'''\cos\tfrac12(a-e)-B'''}{\sin\tfrac12(a-e)},\\[5pt]
B''''=\frac{B'''\cos\tfrac12(b-f)-C'''}{\sin\tfrac12(b-f)}\quad\text{etc.}
\end{gathered}\]
\[\begin{gathered}
A^{\mathrm{V}}=\frac{A''''\cos\tfrac12(a-e)+A'''\sin\tfrac12(a-e)-B''''}{\sin\tfrac12(a-f)},\\[5pt]
B^{\mathrm{V}}=\frac{B''''\cos\tfrac12(b-f)+B'''\sin\tfrac12(b-f)-C''''}{\sin\tfrac12(b-g)}\quad\text{etc.}
\end{gathered}\]
\[
A^{\mathrm{VI}}=\frac{A^{\mathrm{V}}\cos\tfrac12(a-g)-B^{\mathrm{V}}}{\sin\tfrac12(a-g)},\qquad
B^{\mathrm{VI}}=\frac{B^{\mathrm{V}}\cos\tfrac12(b-h)-C^{\mathrm{V}}}{\sin\tfrac12(a-h)}\quad\text{etc.}
\]
etc.

Lex formationis hic satis obvia est, si modo observetur, numeratores in valoribus pro $A''$, $A'''$, $A''''$, $A^{\mathrm{V}}$, $A^{\mathrm{VI}}$ etc. (valor pro $A'$ ab hac regula excipiendus est) alternis vicibus e duabus vel tribus partibus constare.

\InterpolationArticle{13}
\textsc{Theorema.}\quad\textit{Si $X$ est functio arcus $x$ formae $(F)$}
\[
\begin{aligned}
&\alpha+\alpha'\cos x+\alpha''\cos 2x+\alpha'''\cos 3x+\text{etc.}\\[5pt]
&\phantom{\alpha}+\beta'\sin x+\beta''\sin 2x+\beta'''\sin 3x+\text{etc.}
\end{aligned}
\]
\textit{vel huius formae $(G)$}
\[
\begin{aligned}
&\gamma\cos\tfrac12x+\gamma'\cos\tfrac32x+\gamma''\cos\tfrac52x+\text{etc.}\\[5pt]
&+\delta\sin\tfrac12x+\delta'\sin\tfrac32x+\delta''\sin\tfrac52x+\text{etc.}
\end{aligned}
\]
\textit{positoque $x=a$, valor functionis $X$ fit $=0$: erit $X$ divisibilis per $\sin\tfrac12(x-a)$, quotiensque in casu priore formae $G$, in posteriore formae $F$.}

\textit{Demonstratio.\quad Casus prior.}\quad Si in functione $X$ pro quavis parte $\cos nx$ substituitur $\cos nx-\cos na$, pro quavis parte $\sin nx$ autem $\sin nx-\sin na$,
\SourceLeaf{296}
denique pro $\alpha$, $\alpha-\alpha$ sive $0$, manifestum est, partes determinatas hoc modo adiunctas, mutuo se tollere adeoque functionem $X$ hoc modo non variari. Singulae vero partes ipsius $X$ nunc per $\sin\tfrac12(x-a)$ divisibiles erunt, puta
\[
\begin{aligned}
\frac{\cos nx-\cos na}{\sin\tfrac12(x-a)}
&=-2\sin((n-\tfrac12)x+\tfrac12a)\\[5pt]
&\phantom{=}{-}2\sin((n-\tfrac32)x+\tfrac32a)\\[5pt]
&\phantom{=}{-}2\sin((n-\tfrac52)x+\tfrac52a)\\[5pt]
&\phantom{=}{-}\text{etc.}\\[5pt]
&\phantom{=}{-}2\sin((\tfrac12x+(n-\tfrac12)a)
\end{aligned}
\]
\[
\begin{aligned}
\frac{\sin nx-\sin na}{\sin\tfrac12(x-a)}
&=+2\cos((n-\tfrac12)x+\tfrac12a)\\[5pt]
&\phantom{=}{+}2\cos((n-\tfrac32)x+\tfrac32a)\\[5pt]
&\phantom{=}{+}2\cos((n-\tfrac52)x-\tfrac52a)\\[5pt]
&\phantom{=}{+}\text{etc.}\\[5pt]
&\phantom{=}{+}2\cos(\tfrac12x+(n-\tfrac12)a)
\end{aligned}
\]
Uterque coëfficiens manifesto ad formam $G$ reduci potest, quamobrem etiam tota $X$ per $\sin\tfrac12(x-a)$ divisibilis, quotiensque ad formam $G$ reducibilis est. Q.~E.~D.

\textit{Casus posterior.}\quad Si in functione $X$ pro quavis parte $\cos nx$ (designante iam $n$ non ut ante integrum, sed integri imparis semissem) substituitur $\cos nx-\cos\tfrac12(x-a)\cos na$, pro quavis parte $\sin nx$ autem $\sin nx-\cos\tfrac12(x-a)\sin na$, manifestum est, partes, quae hoc modo functioni $X$ accedunt, mutuo se tollere, adeoque $X$ non variari. Singulae autem partes ipsius $X$ nunc per $\sin\tfrac12(x-a)$ divisibiles erunt, puta
\[
\begin{aligned}
\frac{\cos nx-\cos\tfrac12(x-a)\cos na}{\sin\tfrac12(x-a)}\qquad\\[5pt]&=-2\sin((n-\tfrac12)x+\tfrac12a)\\[5pt]
&\phantom{=}{-}2\sin((n-\tfrac32)x+\tfrac32a)\\[5pt]
&\phantom{=}{-}2\sin((n-\tfrac52)x+\tfrac52a)\\[5pt]
&\phantom{=}{-}\text{etc.}\\[5pt]
&\phantom{=}{-}2\sin(x+(n-1)a)\\[5pt]
&\phantom{=}{-}\sin na
\end{aligned}
\]
\[
\begin{aligned}
\frac{\sin nx-\cos\tfrac12(x-a)\sin na}{\sin\tfrac12(x-a)}\qquad\\[5pt]&=+2\cos((n-\tfrac12)x+\tfrac12a)\\[5pt]
&\phantom{=}{+}2\cos((n-\tfrac32)x+\tfrac32a)\\[5pt]
&\phantom{=}{+}2\cos((n-\tfrac52)x+\tfrac52a)\\[5pt]
&\phantom{=}{+}\text{etc.}\\[5pt]
&\phantom{=}{+}2\cos(x+(n-1)a)\\[5pt]
&\phantom{=}{+}\cos na
\end{aligned}
\]
\SourceLeaf{297}
ut per multiplicationem facile confirmatur. Uterque quotiens est formae $F$. Quamobrem functio tota $X$ per $\sin\tfrac12(x-a)$ divisibilis, quotiensque ad formam $F$ reducibilis erit. Q.~E.~D.

\InterpolationArticle{14}
\textsc{Theorema.}\quad\textit{Si functio $X$ formae $F$ vel $G$ pro pluribus valoribus diversis\HistoricalNote{24} ipsius $x$, puta $a$, $b$, $c$, $d$ etc. fit $=0$, per productum}
\[
\sin\tfrac12(x-a)\sin\tfrac12(x-b)\sin\tfrac12(x-c)\sin\tfrac12(x-d)\ .\,.\,.\,.
\]
\textit{divisibilis, quotiensque vel eiusdem formae erit ut $X$ vel diversae, prout multitudo valorum $a$, $b$, $c$, $d$ $.\,.\,.\,.$ par est vel impar.}

\textit{Demonstr.}\quad Ponatur $X=X'\sin\tfrac12(x-a)$, sitque $M$ valor ipsius $X'$ pro $x=b$. Hinc $M\sin\tfrac12(b-a)$ erit valor ipsius $X$ pro $x=b$, qui quum esse debeat $=0$, necessario erit $M=0$. Quare $X'$ divisibilis erit per $\sin\tfrac12(x-b)$ et simili ratione quotiens hinc oriundus per $\sin\tfrac12(x-c)$, et sic porro. Quare $X$ divisibilis erit per productum
\[
\sin\tfrac12(x-a)\sin\tfrac12(x-b)\sin\tfrac12(x-c)\sin\tfrac12(x-d)\ .\,.\,.\quad\text{Q.~E.~D.}
\]
Altera theorematis pars tam obvia est, ut demonstratione opus non sit.

Ceterum aeque obvium est theorema inversum, si $X$ per productum
\[
\sin\tfrac12(x-a)\sin\tfrac12(x-b)\sin\tfrac12(x-c)\sin\tfrac12(x-d)\ .\,.\,.\,.
\]
divisibilis sit, ipsius valorem fieri $=0$, si ipsi $x$ aliquis valorum $a$, $b$, $c$, $d$ etc. tribuatur.

\InterpolationArticle{15}
Ex theoremate praec. sequitur, si functionis $X$, formae $F$, valores pro $x=a$, $b$, $c$, $d$ etc. resp. sint $A$, $B$, $C$, $D$ etc., quamvis aliam similem functionem arcus $x$, quae iisdem valoribus satisfaciat, sub formula $X+PY$ contentam esse debere, ubi per $P$ designamus productum
\[
\sin\tfrac12(x-a)\sin\tfrac12(x-b)\sin\tfrac12(x-c)\sin\tfrac12(x-d)\quad\text{etc.}
\]
per $Y$ autem functionem indefinitam arcus $x$, quae sit vel formae $F$ vel formae
\SourceLeaf{298}
\(G\), prout multitudo valorum \(a,b,c,d\) etc., quam per \(\mu\) designabimus, par est vel impar. Iam observamus:

I. In artt. 10, 11, 12 eruere docuimus functionem \(X\) ordinis \(m^{\mathrm{ti}}\) (i.e. ultra \(\cos mx\) et \(\sin mx\) non progredientem), quae \(\mu\) valoribus propositis \(A,B,C,D\) etc. satisfaciat, ita ut \(m\) sit \(=\frac{1}{2}\mu-\frac{1}{2}\), vel \(\frac{1}{2}\mu\), prout \(\mu\) impar est vel par. Productum \(P\) manifesto est ordinis \({\frac{1}{2}\mu}^{\mathrm{ti}}\). Quando itaque \(\mu\) impar est, adeoque \(Y\) vel ordinis \({\frac{1}{2}}^{\mathrm{ti}}\) vel altioris, erit \(PY\) ordinis ad minimum \({\frac{1}{2}\mu+\frac{1}{2}}^{\mathrm{ti}}\) adeoque ordinis altioris quam \(X\), unde colligitur, functiones alias quam \(X\), quae iisdem valoribus propositis satisfaciant, non dari, nisi ordinis altioris. Quando vero \(\mu\) est par, \(PY\) tunc tantummodo fit ordinis \({\frac{1}{2}\mu}^{\mathrm{ti}}\), si pro \(Y\) accipitur quantitas definita; quare in hoc casu infinite multae quidem functiones similis formae et ordinis ut \(X\) dantur, quae iisdem valoribus satisfaciunt, omnes autem sub forma \(X+hP\) comprehensae erunt, designante \(h\) quantitatem definitam. Eaedem conclusiones ex methodo, per quam in artt. 10, 11 functionem \(X\) derivavimus, sponte sequuntur.

II. Vice versa autem, si quae functio \(X'\) aliunde innotuisset, quae omnibus quidem valoribus propositis \(A,B,C,D\) etc. satisfacit, sed ad altiorem quam opus est ordinem ascendit, functionem \(Y\) ita determinare licebit, ut \(X'+PY\) ad ordinem \({\frac{1}{2}\mu-\frac{1}{2}}^{\mathrm{tum}}\) vel \({\frac{1}{2}\mu}^{\mathrm{tum}}\) (prout \(\mu\) impar est vel par) deprimatur. Sint termini summi in \(X'\), \(K\cos nx+L\sin nx\); in \(P\), \(k\cos\frac{1}{2}\mu x+l\sin\frac{1}{2}\mu x\), accipianturque pro terminis summis in \(Y\) hi
\[
\rho=-\frac{2(Kk+Ll)}{kk+ll}\cos(n-\tfrac{1}{2}\mu)x
     +\frac{2(Kl-Lk)}{kk+ll}\sin(n-\tfrac{1}{2}\mu)x
\]
Hinc calculo facto invenietur, coëfficientes ipsorum \(\cos nx\) et \(\sin nx\) in \(X'+PY\) fieri \(=0\), siquidem fuerit \(n>\frac{1}{2}\mu\); similique modo ex terminis summis functionis evolutae \(X'+P\rho\) (qui erunt ordinis \({n-1}^{\mathrm{ti}}\)) definientur termini proxime inferiores in \(Y\) (ordinis \({n-\frac{1}{2}\mu-1}^{\mathrm{ti}}\)). Haec operatio eousque repetenda erit, donec \(X'+PY\) depressa sit ad ordinem non altiorem quam \(\frac{1}{2}\mu\), adeoque ad ordinem \({\frac{1}{2}\mu-\frac{1}{2}}^{\mathrm{tum}}\), quando \(\mu\) impar est, ad ordinem \({\frac{1}{2}\mu}^{\mathrm{tum}}\), quando \(\mu\) par est. Ulterius hanc depressionem non patere, nullo negotio perspicitur.

III. In casu itaque priore (quando \(\mu\) impar est) hoc modo necessario ad eandem functionem \(X\) perveniemus, quam per methodum supra traditam e valoribus datis \(A,B,C,D\) etc. eruissemus. In casu posteriore autem methodus modo exposita suppeditat functionem, quae quidem eiusdem ordinis erit, cuius est functio, per methodum art. 12, sive per methodum art. 11, si terminus ultimus
\SourceLeaf{299}
\(=0\) ponitur, oriunda, attamen ab hac diversa esse potest. Sit functio illa \(X''\), haec \(X'''\); adeoque \(X'''=X''+hP\), designante \(h\) quantitatem definitam. Ponamus terminos summos in \(X''\) esse \(K'\cos\frac{1}{2}\mu x+L'\sin\frac{1}{2}\mu x\); in \(X'''\) autem \(K''\cos\frac{1}{2}\mu x+L''\sin\frac{1}{2}\mu x\), ita ut sit \(K''=K'+hk\), \(L''=L'+hl\). Iam non difficile est, ex art. 11 vel ex art. 12 demonstrare, si aggregatum \(a+b+c+d+\text{etc.}\) designetur per \(s\), esse \(K''\cos\frac{1}{2}s+L''\sin\frac{1}{2}s=0\), unde deducitur
\[
\begin{aligned}
h&=-\frac{K'\cos\frac{1}{2}s+L'\sin\frac{1}{2}s}
              {k\cos\frac{1}{2}s+l\sin\frac{1}{2}s}\\[5pt]
K''&=\phantom{-}\frac{K'l-L'k}
              {k\cos\frac{1}{2}s+l\sin\frac{1}{2}s}\sin\tfrac{1}{2}s\\[5pt]
L''&=-\frac{K'l-L'k}
              {k\cos\frac{1}{2}s+l\sin\frac{1}{2}s}\cos\tfrac{1}{2}s
\end{aligned}
\]
Haud difficilius demonstratur, esse
\[
k=\pm\frac{\cos\frac{1}{2}s}{2^{\mu-1}},\quad
l=\pm\frac{\sin\frac{1}{2}s}{2^{\mu-1}}
\]
ubi signum superius valet, quando \(\frac{1}{2}\mu\) est par, inferius, quando est impar. Quare in valoribus modo traditis pro \(K''\) et \(L''\) denominator fit \(=\pm\frac{1}{2^{\mu-1}}\) adeoque
\[
\begin{aligned}
K''&=\phantom{-}(K'\sin\tfrac{1}{2}s-L'\cos\tfrac{1}{2}s)\sin\tfrac{1}{2}s\\[5pt]
L''&=-(K'\sin\tfrac{1}{2}s-L'\cos\tfrac{1}{2}s)\cos\tfrac{1}{2}s
\end{aligned}
\]

\InterpolationArticle{16}

Hactenus tales functiones consideravimus, in quibus tum cosinus tum sinus adsunt: saepissime vero aut hi aut illi absunt, quae functionum genera seorsim tractare conveniet. Utrumque quidem casum ad casum generalem hucusque consideratum reducere liceret; attamen etiam magis e re esse videtur, hanc disquisitionem ad problema art. 3 reducere.

\textit{Primo}, si constat, omnes coëfficientes \(\beta',\beta'',\beta'''\) etc. in functione \(X\) (art. 10) esse \(=0\), multitudo incognitarum ad \(m+1\) diminuitur, et proin pro tot valoribus diversis ipsius \(x\) valores respondentes functionis \(X\) novisse sufficit. Quoniam vero in hoc casu valor functionis \(X\) non mutatur, si pro \(x\) substituitur \(-x\), manifestum est, non solum tales valores ipsius \(x\), quorum differentia peripheriae vel multiplo peripheriae aequalis est, sed tales quoque, quorum summa est \(0\) vel peripheriae vel multiplo peripheriae aequalis, pro diversis non esse habendos. Aut generaliter, ut duos valores ipsius \(x\) pro diversis habere liceat, cosinus inaequa\SourceLeaf{300}%
les habere debent. Quibus ita intellectis sint \(m+1\) valores dati functionis \(X\)
hi \(A, B, C, D\) etc., valoribus ipsius \(x\) his \(a, b, c, d\) etc. respondentes, quaeriturque expressio generalis pro valore \(T\), quem functio pro \(x=0\) adipiscitur.
Quum functionem \(X\) in casu nostro etiam sub formam
\[
\gamma+\gamma'\cos x+\gamma''(\cos x)^2+\gamma'''(\cos x)^3+\text{etc.}+\gamma^m(\cos x)^m
\]
reducere liceat, habebimus per art. 3
\[
\begin{aligned}
T={}&\frac{(\cos t-\cos b)(\cos t-\cos c)(\cos t-\cos d)\mathinner{.\,.\,.}}{(\cos a-\cos b)(\cos a-\cos c)(\cos a-\cos d)\mathinner{.\,.\,.}}\,A\\[5pt]
&+\frac{(\cos t-\cos a)(\cos t-\cos c)(\cos t-\cos d)\mathinner{.\,.\,.}}{(\cos b-\cos a)(\cos b-\cos c)(\cos b-\cos d)\mathinner{.\,.\,.}}\,B\\[5pt]
&+\frac{(\cos t-\cos a)(\cos t-\cos b)(\cos t-\cos d)\mathinner{.\,.\,.}}{(\cos c-\cos a)(\cos c-\cos b)(\cos c-\cos d)\mathinner{.\,.\,.}}\,C\\[5pt]
&+\frac{(\cos t-\cos a)(\cos t-\cos b)(\cos t-\cos c)\mathinner{.\,.\,.}}{(\cos d-\cos a)(\cos d-\cos b)(\cos d-\cos c)\mathinner{.\,.\,.}}\,D\\[5pt]
&+\text{etc.}
\end{aligned}
\]

\textit{Secundo}, si omnes coëfficientes \(\alpha,\alpha',\alpha'',\alpha'''\) etc. evanescunt, multitudo incognitarum erit \(m\), sufficitque adeo, valores functionis \(X\) pro totidem valoribus
diversis ipsius \(x\) novisse. Quum valor functionis \(X\) pro \(x=-t\) ex valore pro
\(x=+t\) sponte sequatur (per solam mutationem signi in oppositum), etiam hic
tales valores pro diversis haberi nequeunt, unde facile colligitur, hic perinde ut
in casu praec. eos tantummodo valores ipsius \(x\) pro diversis agnosci, quorum cosinus sunt inaequales. Praeterea quum valor functionis \(X\) pro \(x=0,180^0,360^0\)
etc. et generaliter pro tali valore ipsius \(x\), cuius sinus est \(=0\), sponte fiat \(=0\),
adeoque iam ex natura problematis datus sit, talem valorem inter \(m\) datos non
numerari supponimus. Iam quum constet, esse
\[
\frac{\sin nx}{\sin x}=2\cos(n-1)x+2\cos(n-3)x+2\cos(n-5)x+\mathinner{.\,.\,.}+2\cos x
\]
pro valore pari ipsius \(n\), atque
\[
=2\cos(n-1)x+2\cos(n-3)x+2\cos(n-5)x+\mathinner{.\,.}+2\cos 2x+1
\]
pro valore impari ipsius \(n\), facile concluditur \(\frac{X}{\sin x}\) ad formam
\[
\delta+\delta'\cos x+\delta''(\cos x)^2+\delta'''(\cos x)^3+\text{etc.}+\delta^{m-1}(\cos x)^{m-1}
\]
\SourceLeaf{301}
reduci posse, adeoque, quum valores functionis \(\frac{X}{\sin x}\) pro \(x=a,b,c,d\mathinner{.\,.\,.}\) fiant
\(=\frac{A}{\sin a},\frac{B}{\sin b},\frac{C}{\sin c},\frac{D}{\sin d}\mathinner{.\,.\,.\,.}\), haberi per eundem art. 3
\[
\begin{aligned}
T={}&\frac{\sin t(\cos t-\cos b)(\cos t-\cos c)(\cos t-\cos d)\mathinner{.\,.\,.}}{\sin a(\cos a-\cos b)(\cos a-\cos c)(\cos a-\cos d)\mathinner{.\,.\,.}}\,A\\[5pt]
&+\frac{\sin t(\cos t-\cos a)(\cos t-\cos c)(\cos t-\cos d)\mathinner{.\,.\,.}}{\sin b(\cos b-\cos a)(\cos b-\cos c)(\cos b-\cos d)\mathinner{.\,.\,.}}\,B\\[5pt]
&+\frac{\sin t(\cos t-\cos a)(\cos t-\cos b)(\cos t-\cos d)\mathinner{.\,.\,.}}{\sin c(\cos c-\cos a)(\cos c-\cos b)(\cos c-\cos d)\mathinner{.\,.\,.}}\,C\\[5pt]
&+\frac{\sin t(\cos t-\cos a)(\cos t-\cos b)(\cos t-\cos c)\mathinner{.\,.\,.}}{\sin d(\cos d-\cos a)(\cos d-\cos b)(\cos d-\cos c)\mathinner{.\,.\,.}}\,D\\[5pt]
&+\text{etc.}
\end{aligned}
\]
Ceterum in utroque casu formula pro \(T\) cum functione \(X\) identica erit, mutando
\(t\) in \(x\), quoniam illa indefinite pro valore quocunque ipsius \(t\) valet.

\InterpolationArticle{17}
E principiis algebraicis facile conclusio petitur, si, pro casu priore art. praec.
productum
\[
(\cos t-\cos a)(\cos t-\cos b)(\cos t-\cos c)(\cos t-\cos d)\text{ etc.}
\]
pro posteriore autem productum
\[
\sin t(\cos t-\cos a)(\cos t-\cos b)(\cos t-\cos c)(\cos t-\cos d)\text{ etc.}
\]
per \(P\) designetur, quamvis functionem ipsi \(X\) similem, quae iisdem valoribus
\(A,B,C,D\) etc. pro \(x=a,b,c,d\) etc. satisfacit, necessario sub forma \(X+PY\)
contentam esse debere, ubi \(Y\) est functio indefinita arcus \(x\) formae
\[
\gamma+\gamma'\cos x+\gamma''\cos 2x+\gamma'''\cos 3x+\text{etc.}
\]
Theorema inversum, scilicet quamlibet functionem sub forma \(X+PY\) contentam illis valoribus satisfacere, sponte patet. Porro manifestum est, productum
\(P\) esse ordinis \(m+1^{\mathrm{ti}}\), adeoque proxime altioris quam \(X\); quamobrem alia functio
ipsi \(X\) similis quae iisdem valoribus satisfaciat non datur, nisi ex ordine altiore
quam \(X\).

Vice versa, si quae functio \(X'\), ipsi \(X\) similis quidem, sed altioris ordinis,
aliunde innotuit, quae valoribus datis \(A,B,C,D\) etc. satisfacit. functionem \(Y\),
eius quam tradidimus formae, ita determinare licet, ut \(X'+PY\) ad ordinem
\endgroup
